\section{Study Design}

\subsection{Prerequisite evidence gate}

Before exposing child participants to a dedicated bilateral-audio experiment, the literature review and adult/technical feasibility work should establish that the acoustic pattern is tolerable, that the proposed timing has a plausible non-trauma-specific mechanism, and that the study is worth the additional participant burden.

\subsection{Design}

A randomized within-participant crossover or micro-randomized design is preferred for the initial mechanism study. Eligible challenge-recovery events are randomized among conditions A, B, and C subject to safety and carryover rules. Eligibility, child-requested non-participation and any session-level safety hold are decided and logged before revealing the event assignment. The randomization schedule is generated independently of current task performance except for preregistered eligibility criteria. Assignment, delivered condition, protocol deviation and reason for non-delivery are recorded separately; postassignment refusal does not silently remove an event from the assigned-arm denominator.

\subsection{Acoustic control}

The alternating and non-alternating auditory conditions must be matched on stimulus identity, average level, peak level, spectral content, event count, duration, rhythm and perceived salience as far as a non-alternating pattern permits. Per-ear energy, channel balance and stereo routing must be measured on qualified private phone/headphone configurations, rather than inferred from digital waveform equality; any unavoidable side-distribution difference is documented. Acceptance tolerances, test signals, calibration procedure, equipment eligibility and failure handling must be frozen in a separate acoustic-validation record before participant exposure. The study should not compare a salient alternating sound with an acoustically weaker or less engaging control.

\subsection{Outcomes}

The primary proximal outcome is recovery latency. Secondary outcomes can include unrecovered-event rate, post-challenge task accuracy, child-rated comfort, perceived distraction, subjective activation, and voluntary stop requests.

ADHD symptom change is not a primary outcome of this proof-of-mechanism study. Any longer-term clinical claim would require a separate adequately powered trial after the proximal mechanism is established.

\subsection{Carryover and dose}

Inter-event washout, maximum number of experimental exposures per session, maximum sound level, and immediate stop criteria must be fixed before data collection and tested in technical/adult feasibility work. The stop control must remain available throughout exposure; an immediate child stop or safety override ends sound delivery and is recorded with its reason and timing. A subsequent event is eligible only after the prespecified washout and renewed safety/comfort check. Device failure, protocol deviation, discomfort and missing telemetry have distinct codes and analysis roles. Events randomized before a stop remain in the assigned-arm analysis as specified in the analysis section; exclusions are limited to preassignment ineligibility and cannot depend on observed recovery.
